\chapter{常用不等式与估计方法索引：习题解答}

\begin{solution}{C.1}
对 \(n\) 归纳：\(n=2\) 是三角不等式；若对 \(n\) 成立，则
\[
 \left|\sum_{k=1}^{n+1}x_k\right|
 \le\left|\sum_{k=1}^{n}x_k\right|+|x_{n+1}|
 \le\sum_{k=1}^{n+1}|x_k|.
\]
范数情形由 \(x=(x-y)+y\) 得 \(\|x\|-\|y\|\le\|x-y\|\)，交换 \(x,y\) 后合并。
\end{solution}

\begin{solution}{C.2}
若 \(y=0\) 结论直接成立。对实数 \(t\)，
\[
 0\le\|x-ty\|^2=\|x\|^2-2t\langle x,y\rangle+t^2\|y\|^2.
\]
该二次多项式判别式不大于零，故
\[
 |\langle x,y\rangle|^2\le\|x\|^2\|y\|^2.
\]
等号当且仅当二次式对某个 \(t\) 为零，即 \(x=ty\)，也就是线性相关。
\end{solution}

\begin{solution}{C.3}
令 \(X=(\sum|x_k|^p)^{1/p}\)、\(Y=(\sum|y_k|^q)^{1/q}\)。若 \(XY=0\) 结论直接成立。对 \(u_k=|x_k|/X,\ v_k=|y_k|/Y\)，Young 不等式给出
\[
 u_kv_k\le\frac{u_k^p}{p}+\frac{v_k^q}{q}.
\]
求和得 \(\sum u_kv_k\le1/p+1/q=1\)，再乘 \(XY\)。
\end{solution}

\begin{solution}{C.4}
共轭指数 \(q=p/(p-1)\)，故 \((p-1)q=p\)。若 \(\|x+y\|_p=0\)，结论直接成立。否则
\[
 \|x+y\|_p^p
 \le(\|x\|_p+\|y\|_p)
 \left(\sum|x+y|^{(p-1)q}\right)^{1/q}
 =(\|x\|_p+\|y\|_p)\|x+y\|_p^{p-1}.
\]
除以最后因子即得结论。
\end{solution}

\begin{solution}{C.5}
先由凸性定义证明两点权重情形。对 \(n\) 归纳：令 \(\mu=1-\lambda_n\)。若 \(\mu>0\)，把前 \(n-1\) 项归一化为权重 \(\lambda_i/\mu\)，先应用归纳假设，再与 \(x_n\) 应用两点凸性。若 \(\mu=0\) 结论为等式。严格凸时，每次正权两点取等都要求两点相同，故所有正权对应的 \(x_i\) 相等；反向显然。
\end{solution}

\begin{solution}{C.6}
令 \(v_n=A+\sum_{k<n}b_ku_k\)。则 \(u_n\le v_n\) 且
\[
 v_{n+1}\le(1+b_n)v_n.
\]
归纳得到 \(u_n\le A\prod_{k<n}(1+b_k)\)，再用 \(1+t\le e^t\)。

若 \(A_n\) 非负单调递增，固定 \(n\) 并把所有较早的 \(A_k\) 以 \(A_n\) 控制，同一论证给出
\[
 u_n\le A_n\prod_{k<n}(1+b_k).
\]
\end{solution}

\begin{solution}{C.7}
对任意 \(\eta>0\)，原不等式蕴含
\[
 u(t)\le \eta+\int_0^tb(s)u(s)\,ds.
\]
按 \(A=\eta\) 的已证情形，
\[
 u(t)\le\eta\exp\left(\int_0^tb(s)\,ds\right).
\]
右端对固定 \(t\) 随 \(\eta\downarrow0\) 趋零，且 \(u(t)\ge0\)，故 \(u(t)=0\)（在原不等式成立的点上）。
\end{solution}

\begin{solution}{C.8}
分情形讨论 \(x,y\) 与区间 \([-M,M]\) 的位置。若都在区间内，差不变；若都在同一外侧，差为零；若一内一外，投影到端点只会缩短距离；若分居两侧，截断后距离 \(2M\le|x-y|\)。因此
\[
 |T_M(x)-T_M(y)|\le|x-y|.
\]
\end{solution}

\begin{solution}{C.9}
若三项误差分别严格小于 \(\varepsilon/2,\varepsilon/4,\varepsilon/4\)，三角不等式给出总误差严格小于 \(\varepsilon\)。与平均分配 \(\varepsilon/3\) 相同，关键是各份为正且总和不超过目标，并按依赖关系依次选择辅助对象和指标。
\end{solution}

\begin{solution}{C.10}
令 \(s_n=n^{-1}\sum_{k=1}^na_k\)。给定 \(\varepsilon>0\)，取 \(N\) 使 \(k\ge N\) 时 \(|a_k-L|<\varepsilon/2\)。固定常数 \(C=\sum_{k<N}|a_k-L|\)，取 \(n\) 大到 \(C/n<\varepsilon/2\)，则
\[
 |s_n-L|\le C/n+\varepsilon/2<\varepsilon.
\]
反例 \(a_n=(-1)^n\) 不收敛，而其平均的绝对值不超过 \(1/n\)，故趋于 \(0\)。
\end{solution}

\begin{solution}{C.11}
有限极限版本按正文：从某项起差商夹在 \(L\pm\varepsilon\) 之间，乘正差 \(b_{k+1}-b_k\) 后求和，除以 \(b_n\)，固定前缀项消失。

若差商趋于 \(+\infty\)，给定 \(M>0\)，从某项起
\[
 a_{k+1}-a_k\ge M(b_{k+1}-b_k).
\]
求和得
\[
 \frac{a_n}{b_n}\ge
 M+\frac{a_N-Mb_N}{b_n}.
\]
令 \(n\to\infty\) 得下极限至少为 \(M\)。因 \(M\) 任意，\(a_n/b_n\to+\infty\)。
\end{solution}

\begin{solution}{C.12}
取 \(a_n=H_n\)、\(b_n=\log n\)。分母严格递增且趋于无穷。差商为
\[
 \frac{H_{n+1}-H_n}{\log(n+1)-\log n}
 =\frac{1/(n+1)}{\log(1+1/n)}.
\]
由 \(\log(1+t)\sim t\)，该比值趋于 \(1\)。Stolz--Cesàro 定理给出 \(H_n/\log n\to1\)。
\end{solution}
